\chapter{Variables and the Idea of Phase}\label{ch:variables-phase}

\index{variables}\index{phase system}

In most languages, declaring a variable is a mundane act.
You pick a name, assign a value, and move on.
In Lattice, declaring a variable is a \emph{decision}---a declaration of intent about how that value should behave over its lifetime.
Will it change?
Should it be locked down?
Or should the language figure that out for you?

Lattice gives you three keywords for variable declaration---\lstinline|let|, \lstinline|flux|, and \lstinline|fix|---and each one makes a different promise to the runtime.
Behind these keywords lies the \emph{phase system}\index{phase system}, the feature that gives Lattice its name and its crystalline metaphor.
This chapter introduces both the practical mechanics of variable binding and the philosophy that makes them meaningful.

%% ───────────────────────────────────────────────────────────
\section{\texttt{let}, \texttt{flux}, and \texttt{fix}}\label{sec:let-flux-fix}

\subsection{\texttt{let} --- Inferred Phase}\label{subsec:let}
\index{let}

The \lstinline|let| keyword declares a variable and infers its phase from the value you assign:

\begin{lstlisting}[language=Lattice, caption={Declaring variables with let}]
let name = "Lattice"      // a string
let version = 28           // an integer
let pi = 3.14159           // a float
let active = true          // a boolean
let colors = ["red", "green", "blue"]  // an array
\end{lstlisting}

With \lstinline|let|, the variable's mutability depends on what you assign.
In casual mode (the default), \lstinline|let| bindings are treated as fluid---you \emph{can} reassign them or mutate their contents:

\begin{lstlisting}[language=Lattice, caption={let bindings are fluid in casual mode}]
let score = 100
// In casual mode, this works:
// score = 200

let items = [1, 2, 3]
items.push(4)  // modifying in place works
print(items)   // [1, 2, 3, 4]
\end{lstlisting}

Think of \lstinline|let| as the ``just give me a variable'' keyword---quick, convenient, and unassuming.
It's perfect for scripts, quick prototypes, and situations where you don't need to think about mutability.

\subsection{\texttt{flux} --- Explicitly Fluid}\label{subsec:flux}
\index{flux}\index{fluid phase}

The \lstinline|flux| keyword (pronounced like ``flux'' in metallurgy) declares a variable that is \emph{explicitly fluid}\index{fluid phase}.
A fluid value can be modified, reassigned, grown, and reshaped:

\begin{lstlisting}[language=Lattice, caption={Declaring mutable variables with flux}]
flux counter = 0
counter += 1
counter += 1
print(counter)  // 2

flux inventory = Map::new()
inventory["apples"] = 12
inventory["bananas"] = 6
print(inventory)  // {"apples": 12, "bananas": 6}

flux temperatures = []
temperatures.push(72.1)
temperatures.push(68.5)
temperatures.push(74.3)
print(temperatures)  // [72.1, 68.5, 74.3]
\end{lstlisting}

The \lstinline|flux| prefix is a signal to anyone reading your code: ``this variable \emph{will} change.''
When you see \lstinline|flux| in a program, you know to watch for mutations.

\begin{defnbox}{Fluid Phase}
A value in the \emph{fluid phase}\index{fluid phase} is mutable.
Arrays can be pushed to, maps can gain new keys, and the variable can be reassigned.
In the chemistry metaphor, a fluid value is like molten metal---it has not yet solidified into its final form.
You declare fluid variables with \lstinline|flux| (or the shorthand prefix \lstinline|~|).
\end{defnbox}

Under the hood, the compiler emits an \lstinline|OP_MARK_FLUID|\index{OP\_MARK\_FLUID} instruction when it encounters \lstinline|flux|, which sets the value's phase tag to \lstinline|VTAG_FLUID| (you can see this in \filename{src/stackcompiler.c}).
This tag travels with the value for its entire lifetime.

\subsection{\texttt{fix} --- Explicitly Crystal}\label{subsec:fix}
\index{fix}\index{crystal phase}

The \lstinline|fix| keyword declares a variable that is \emph{crystal}\index{crystal phase}---immutable, locked, permanent.
The value must already be frozen (or will be frozen by the declaration):

\begin{lstlisting}[language=Lattice, caption={Declaring immutable variables with fix}]
fix pi = freeze(3.14159)
fix greeting = freeze("Hello, World!")
fix primes = freeze([2, 3, 5, 7, 11])

print(pi)        // 3.14159
print(greeting)  // Hello, World!
print(primes)    // [2, 3, 5, 7, 11]
\end{lstlisting}

Attempting to modify a crystal value produces an error:

\begin{lstlisting}[language=Lattice, caption={Crystal values reject modification}]
fix config = freeze([1, 2, 3])
// config.push(4)  -- error: cannot mutate crystal value
// config = [4, 5, 6]  -- error: cannot assign to crystal binding
\end{lstlisting}

\begin{defnbox}{Crystal Phase}
A value in the \emph{crystal phase}\index{crystal phase} is immutable.
Its contents cannot be modified, its structure cannot change, and it can be safely shared across threads.
In the chemistry metaphor, a crystal is solid, ordered, and permanent---like a diamond formed under pressure.
You declare crystal variables with \lstinline|fix| (or the shorthand prefix \lstinline|*|).
\end{defnbox}

When the compiler encounters \lstinline|fix|, it emits an \lstinline|OP_FREEZE|\index{OP\_FREEZE} instruction that transitions the value to the crystal phase.

\subsection{The Shorthand Prefixes: \texttt{\~{}} and \texttt{*}}\label{subsec:shorthand-prefixes}
\index{~@\texttt{\~{}} (flux shorthand)}\index{*@\texttt{*} (fix shorthand)}

In type annotations and function parameters, you'll sometimes see \lstinline|~| for fluid and \lstinline|*| for crystal:

\begin{lstlisting}[language=Lattice, caption={Phase shorthand in function parameters}]
fn mutate(data: ~Map) {
    data.set("modified", true)
}

fn inspect(data: *Map) {
    print(data.get("name"))
}
\end{lstlisting}

Here, \lstinline|~Map| means ``a Map that must be fluid'' and \lstinline|*Map| means ``a Map that must be crystal.''
The runtime checks these constraints on every call---pass a crystal map to \lstinline|mutate()| and you'll get a clear error.

These shorthand prefixes are equivalent to writing \lstinline|flux Map| and \lstinline|fix Map| in type position.
Use whichever reads better to you.

\subsection{Comparing the Three Keywords}\label{subsec:comparing-keywords}

\begin{table}[h]
\centering
\begin{tabular}{llll}
\toprule
\textbf{Keyword} & \textbf{Phase} & \textbf{Mutable?} & \textbf{When to Use} \\
\midrule
\lstinline|let|  & Inferred   & Depends on mode & Scripts, quick code, ``just works'' \\
\lstinline|flux| & Fluid      & Yes             & Values you intend to modify \\
\lstinline|fix|  & Crystal    & No              & Constants, shared data, configs \\
\bottomrule
\end{tabular}
\caption{Variable declaration keywords}\label{tab:declaration-keywords}
\end{table}

A rule of thumb: if you're writing a script or prototyping, \lstinline|let| is fine.
If you're writing production code or library code, prefer \lstinline|flux| and \lstinline|fix| for clarity.
In strict mode, \lstinline|let| is actually \emph{disallowed}---you must choose explicitly.

%% ───────────────────────────────────────────────────────────
\section{Why Mutability Is a First-Class Concept}\label{sec:mutability-first-class}
\index{mutability!as first-class concept}

Most languages treat mutability as a binary toggle.
In Rust, you have \lstinline|let| vs.\ \lstinline|let mut|.
In JavaScript, you have \lstinline|const| vs.\ \lstinline|let|.
In Python, everything is mutable and you just hope for the best.

Lattice goes further.
Mutability is not just a property of a \emph{variable}---it's a property of the \emph{value itself}.
This distinction is subtle but profound.

\begin{lstlisting}[language=Lattice, caption={Phase is a property of the value, not the variable}]
flux items = [1, 2, 3]
print(phase_of(items))   // fluid

freeze(items)
print(phase_of(items))   // crystal

// The variable `items` still exists.
// But the *value* it holds has changed phase.
// items.push(4)  -- error: the value is crystal
\end{lstlisting}

The phase travels with the value.
If you pass a frozen array to a function, that function cannot mutate it---even if it assigns it to a new \lstinline|flux| variable.
The crystal phase is embedded in the value, not in the binding.

\begin{lstlisting}[language=Lattice, caption={Phase follows the value through function calls}]
fn try_to_modify(data: any) {
    flux local = data
    // local.push(99)  -- error if data is crystal!
    print(phase_of(local))
}

flux numbers = [1, 2, 3]
try_to_modify(numbers)    // prints: fluid

freeze(numbers)
try_to_modify(numbers)    // prints: crystal (push would fail)
\end{lstlisting}

This design has a powerful consequence: when a value is frozen, it is safe to share it across threads, pass it to unknown functions, or store it in a long-lived data structure---because \emph{nothing} can modify it.
The guarantee is enforced by the runtime, not by programmer discipline.

\begin{notebox}{The Value vs.\ the Binding}
The \lstinline|flux|/\lstinline|fix| keywords control two things:
\begin{enumerate}
  \item Whether the \textbf{variable binding} can be reassigned (pointing to a different value).
  \item The initial \textbf{phase of the value} (fluid or crystal).
\end{enumerate}
The second property persists even after the value is passed to other functions or stored in data structures.
The first property only applies to the original variable name in its scope.
\end{notebox}

%% ───────────────────────────────────────────────────────────
\section{Quick Introduction to \texttt{freeze()} and \texttt{thaw()}}\label{sec:freeze-thaw}
\index{freeze}\index{thaw}

We've already seen \lstinline|freeze()| in passing.
Let's give it proper attention.

\subsection{\texttt{freeze()} --- Fluid to Crystal}\label{subsec:freeze-detail}

\lstinline|freeze()| transitions a value from the fluid phase to the crystal phase.
Once frozen, the value is immutable:

\begin{lstlisting}[language=Lattice, caption={Freezing a value}]
flux settings = Map::new()
settings["debug"] = false
settings["verbose"] = true
settings["max_retries"] = 3

print(phase_of(settings))  // fluid
freeze(settings)
print(phase_of(settings))  // crystal

// settings["debug"] = true  -- error: cannot mutate crystal value
\end{lstlisting}

For collections (arrays, maps, sets), \lstinline|freeze()| is \textbf{deep}\index{freeze!deep}: it freezes the container and all values inside it recursively.

\begin{lstlisting}[language=Lattice, caption={Deep freeze}]
flux data = [[1, 2], [3, 4]]
freeze(data)

// data.push([5, 6])  -- error: outer array is crystal
// data[0].push(3)    -- error: inner array is also crystal
\end{lstlisting}

\subsection{\texttt{thaw()} --- Crystal to Fluid}\label{subsec:thaw-detail}

\lstinline|thaw()| creates a \textbf{mutable copy}\index{thaw!creates copy} of a crystal value.
The original remains frozen:

\begin{lstlisting}[language=Lattice, caption={Thawing a frozen value}]
fix original = freeze([10, 20, 30])
print(phase_of(original))  // crystal

flux copy = thaw(original)
print(phase_of(copy))      // fluid

copy.push(40)
print(copy)      // [10, 20, 30, 40]
print(original)  // [10, 20, 30] (unchanged!)
\end{lstlisting}

The metaphor is precise: thawing a crystal doesn't destroy the crystal.
It melts a copy, leaving the original intact.
This is a key property for safe concurrent programming---multiple threads can thaw the same frozen data independently without interfering with each other.

\subsection{\texttt{clone()} --- Independent Copy}\label{subsec:clone-detail}
\index{clone}

While we're talking about copying, let's mention \lstinline|clone()|.
Unlike \lstinline|thaw()|, which always produces a fluid copy from a crystal, \lstinline|clone()| creates an independent deep copy of any value, preserving its current phase:

\begin{lstlisting}[language=Lattice, caption={Cloning preserves phase}]
flux original = [1, 2, 3]
flux copy = clone(original)

copy.push(4)
print(original)  // [1, 2, 3] (unaffected)
print(copy)      // [1, 2, 3, 4]
\end{lstlisting}

\begin{table}[h]
\centering
\begin{tabular}{lll}
\toprule
\textbf{Function} & \textbf{Input Phase} & \textbf{Output Phase} \\
\midrule
\lstinline|freeze(v)| & Fluid & Crystal (same value, now locked) \\
\lstinline|thaw(v)| & Crystal & Fluid (new copy) \\
\lstinline|clone(v)| & Any & Same as input (new copy) \\
\bottomrule
\end{tabular}
\caption{Phase transition functions}\label{tab:phase-transitions}
\end{table}

We'll explore the full depth of phase transitions---including \lstinline|sublimate|, \lstinline|crystallize|, \lstinline|borrow|, \lstinline|forge|, bonds, reactions, and more---in Part~III of this book.
For now, \lstinline|freeze()|, \lstinline|thaw()|, and \lstinline|clone()| are the essential trio.

%% ───────────────────────────────────────────────────────────
\section{\texttt{\#mode casual} vs.\ \texttt{\#mode strict}}\label{sec:casual-strict}
\index{casual mode}\index{strict mode}\index{\#mode}

Lattice offers two modes of phase enforcement, controlled by a directive at the top of a file:

\begin{lstlisting}[language=Lattice, caption={Mode directives}]
#mode casual   // the default: relaxed phase rules
#mode strict   // tighter enforcement: no let, consume on freeze
\end{lstlisting}

If you don't specify a mode, the default is \lstinline|casual|.

\subsection{Casual Mode}\label{subsec:casual-mode}
\index{casual mode!rules}

In casual mode, Lattice is forgiving:

\begin{itemize}
  \item \lstinline|let| is allowed and produces an inferred-phase binding.
  \item \lstinline|freeze()| transitions a value in place but does \emph{not} consume the variable---you can still read from the binding.
  \item Phase mismatches at function call boundaries produce runtime errors, not compile-time errors.
  \item You can mix \lstinline|let|, \lstinline|flux|, and \lstinline|fix| freely.
\end{itemize}

\begin{lstlisting}[language=Lattice, caption={Casual mode: forgiving}]
let data = [1, 2, 3]
data.push(4)         // works: let is fluid in casual mode
freeze(data)
print(data)          // works: data is still accessible (now crystal)
// data.push(5)      -- error: crystal value
\end{lstlisting}

Casual mode is designed for scripting, exploratory programming, and situations where you want the phase system's safety net without its full strictness.

\subsection{Strict Mode}\label{subsec:strict-mode}
\index{strict mode!rules}

Strict mode tightens the rules:

\begin{itemize}
  \item \lstinline|let| is \textbf{disallowed}\index{strict mode!no let}. You must use \lstinline|flux| or \lstinline|fix| for every declaration, making your intent explicit.
  \item \lstinline|freeze()| \textbf{consumes the binding}\index{strict mode!consume on freeze}. After freezing a variable, the original name is no longer accessible---you must capture the frozen result in a new binding.
  \item The phase checker runs \textbf{before execution}\index{strict mode!static checking}, catching phase errors as compile-time warnings rather than runtime surprises.
  \item Assigning to a crystal binding is a compile-time error.
\end{itemize}

\begin{lstlisting}[language=Lattice, caption={Strict mode: explicit and consuming}]
#mode strict

flux data = [1, 2, 3]
data.push(4)

fix frozen = freeze(data)
// `data` is no longer accessible in strict mode.
// You must use `frozen` from now on.
print(frozen)  // [1, 2, 3, 4]
\end{lstlisting}

The consume-on-freeze behavior prevents a common bug: modifying a value after you've shared a ``frozen'' reference to it.
In casual mode, you could freeze a value, hand the frozen reference to another part of the program, and then accidentally mutate the original.
Strict mode makes this impossible by removing the original name from scope.

\begin{lstlisting}[language=Lattice, caption={Strict mode catches errors early}]
#mode strict

// let x = 42  -- error: use flux or fix instead of let

flux counter = 0
counter += 1

fix pi = freeze(3.14159)
// pi = 2.71828  -- error: cannot assign to crystal binding

// This error is caught before the program runs.
\end{lstlisting}

\begin{tipbox}{When to Use Strict Mode}
Use \lstinline|#mode strict| when:
\begin{itemize}
  \item You're writing a library that will be used by others.
  \item You're writing code that runs in production.
  \item You're working on concurrent code where phase safety is critical.
  \item You want the compiler to catch mistakes before runtime.
\end{itemize}
Use casual mode (the default) when:
\begin{itemize}
  \item You're scripting or prototyping.
  \item You're exploring in the REPL\@.
  \item You want to move fast without ceremony.
\end{itemize}
\end{tipbox}

\subsection{Inside the Phase Checker}\label{subsec:phase-checker-internals}
\index{phase checker}

When you use \lstinline|#mode strict|, the parser sets the program's mode to \lstinline|MODE_STRICT|, which triggers the phase checker (\filename{src/phase\_check.c}) before execution begins.
The phase checker is a static analysis pass that walks the AST and:

\begin{enumerate}
  \item Tracks the phase of every variable through a scope stack.
  \item Verifies that \lstinline|fix| bindings are never assigned to.
  \item Ensures that arguments passed to phase-constrained parameters (\lstinline|~Map|, \lstinline|*Map|) have compatible phases.
  \item Rejects \lstinline|let| bindings---requiring explicit \lstinline|flux| or \lstinline|fix|.
  \item Inside \lstinline|spawn| blocks, checks that fluid bindings from the outer scope are not used across the thread boundary.
\end{enumerate}

The phase checker is defined as a \lstinline|PhaseChecker| struct that maintains a stack of scopes, each mapping variable names to their known phases (\lstinline|PHASE_FLUID|, \lstinline|PHASE_CRYSTAL|, or \lstinline|PHASE_UNSPECIFIED|).
As it walks the AST, it pushes and pops scopes at function and block boundaries, just like the runtime would.

The errors it produces are not fatal---they're warnings that appear before execution.
But they are informative:

\begin{lstlisting}[numbers=none]
phase error: strict mode: use 'flux' or 'fix' instead of 'let' for binding 'x'
phase error: strict mode: cannot assign to crystal binding 'config'
phase error: strict mode: cannot use fluid binding 'data' across thread boundary in spawn
\end{lstlisting}

%% ───────────────────────────────────────────────────────────
\section{First Encounter with the Phase Philosophy}\label{sec:phase-philosophy}
\index{phase system!philosophy}

Now that you've seen the mechanics---\lstinline|flux|, \lstinline|fix|, \lstinline|freeze()|, \lstinline|thaw()|, casual mode, strict mode---let's step back and think about \emph{why} Lattice works this way.

\subsection{The Lifecycle of Data}\label{subsec:data-lifecycle}

Most data in a program goes through a predictable lifecycle:

\begin{enumerate}
  \item \textbf{Construction}: The data is built up incrementally.
     You add items to a list, set fields in a configuration, accumulate results.
  \item \textbf{Stabilization}: At some point, the data is ``done.''
     The configuration is complete, the list has all its items, the computation is finished.
  \item \textbf{Consumption}: The stabilized data is read, shared, serialized, sent to another thread, or stored for later use.
\end{enumerate}

In most languages, there is no formal boundary between these stages.
You just \emph{stop modifying} the data and hope nobody else modifies it later.
This works fine in small programs.
In large programs, with shared state and concurrent threads, it becomes a source of bugs.

Lattice makes the boundary explicit:

\begin{lstlisting}[language=Lattice, caption={The lifecycle of data in Lattice}]
// 1. Construction (fluid)
flux report = Map::new()
report["title"] = "Q4 Revenue"
report["total"] = 1_250_000
report["approved"] = false

// 2. Stabilization (freeze)
report["approved"] = true
freeze(report)

// 3. Consumption (crystal, safe to share)
print(report.get("title"))     // Q4 Revenue
// send_to_archive(report)     // safe: report is immutable
// spawn { process(report) }   // safe: can be sent across threads
\end{lstlisting}

The \lstinline|freeze()| call is the boundary.
Before it, the data is under construction.
After it, the data is permanent.
The runtime enforces this boundary---you can't go back to the construction phase (unless you explicitly \lstinline|thaw()| a copy).

\subsection{The Chemistry Metaphor}\label{subsec:chemistry-metaphor}
\index{phase system!chemistry metaphor}

The terminology in Lattice is deliberately drawn from chemistry and materials science\index{materials science}:

\begin{itemize}
  \item \phase{Fluid} values are like molten metal---shapeable, hot, unstable.
  \item \phase{Crystal} values are like a diamond or quartz---solid, ordered, permanent.
  \item \lstinline|freeze()| is \emph{crystallization}---the transition from fluid to solid.
  \item \lstinline|thaw()| is \emph{melting}---creating a fluid copy from a solid.
  \item \lstinline|forge| blocks are a \emph{forge}---a controlled environment where you heat, shape, and cool a material in a single operation.
\end{itemize}

This metaphor extends further in Part~III, where we'll meet:
\begin{itemize}
  \item \textbf{Sublimation}: shallow freezing that locks the structure but leaves inner values mutable.
  \item \textbf{Bonds}: links between values that cause one to freeze when another does.
  \item \textbf{Reactions}: callbacks that fire when a value changes phase.
  \item \textbf{Pressure}: restrictions on what can change without fully freezing.
  \item \textbf{Alloys}: structs with per-field phase declarations.
  \item \textbf{Arenas}: memory regions where crystal values are stored.
\end{itemize}

But even without these advanced features, the core idea is powerful: \textbf{make the lifecycle of your data visible in the code}.

\subsection{Phase and Concurrency}\label{subsec:phase-concurrency}
\index{phase system!concurrency}

The phase system's biggest payoff comes with concurrency.
Lattice's \lstinline|scope|/\lstinline|spawn| structured concurrency model requires that any value sent across a channel must be \emph{crystal}.
This is enforced at runtime:

\begin{lstlisting}[language=Lattice, caption={Phase and concurrency}]
let ch = Channel::new()

flux data = [1, 2, 3]
// ch.send(data)  -- error: cannot send fluid value

freeze(data)
ch.send(data)  // works: crystal values are safe to share

scope {
    spawn {
        let received = ch.recv()
        print(received)  // [1, 2, 3]
    }
}
\end{lstlisting}

This constraint means that data races---the most common and hardest-to-debug class of concurrency bugs---are impossible in Lattice.
If a value can cross a thread boundary, it is guaranteed to be immutable.
If it is mutable, it cannot cross.
The phase system draws a bright line between ``private, mutable, single-threaded'' and ``shared, immutable, multi-threaded.''

We'll explore concurrency in full in Part~V\@.
For now, the key insight is: the phase system is not just about preventing accidental mutation.
It's a foundation for safe concurrency.

\subsection{A Preview: The Forge Block}\label{subsec:forge-preview}
\index{forge}

One of the most elegant patterns in Lattice is the \lstinline|forge| block, which encapsulates the entire lifecycle of a value---construction, stabilization, and output---in a single expression:

\begin{lstlisting}[language=Lattice, caption={Forge: build fluid, end crystal}]
fix server_config = forge {
    flux temp = Map::new()
    temp.set("host", "0.0.0.0")
    temp.set("port", "443")
    temp.set("tls", "true")
    temp.set("max_connections", "1024")
    freeze(temp)
}

// server_config is crystal, ready to share
print(server_config.get("host"))  // 0.0.0.0
print(phase_of(server_config))    // crystal
\end{lstlisting}

Inside the \lstinline|forge| block, you work with fluid data.
The block's result is automatically treated as crystal.
The temporary mutable state never escapes.
This pattern is so useful that you'll see it throughout Lattice codebases wherever immutable configurations, frozen caches, or thread-safe snapshots are built.

%% ───────────────────────────────────────────────────────────
\section{Practical Patterns}\label{sec:variable-patterns}

Let's close with some practical patterns that show how \lstinline|let|, \lstinline|flux|, \lstinline|fix|, and \lstinline|freeze()| work together in real code.

\subsection{Accumulate Then Freeze}\label{subsec:accumulate-freeze}
\index{patterns!accumulate then freeze}

Build a collection incrementally, then freeze it for safe consumption:

\begin{lstlisting}[language=Lattice, caption={Accumulate then freeze}]
fn load_user_ids(filename: String) -> [Int] {
    let content = read_file(filename)
    let lines = content.split("\n")
    flux ids = []
    for line in lines {
        let trimmed = line.trim()
        if !trimmed.is_empty() {
            ids.push(parse_int(trimmed))
        }
    }
    freeze(ids)
    return ids
}
\end{lstlisting}

The function returns a frozen array.
Callers can trust that the data will not change out from under them.

\subsection{Configuration Objects}\label{subsec:config-objects}
\index{patterns!configuration objects}

Build a configuration with sensible defaults, allow overrides, then lock it down:

\begin{lstlisting}[language=Lattice, caption={Configuration with defaults}]
fn make_config(overrides: Map) -> Map {
    flux cfg = Map::new()

    // Defaults
    cfg["timeout"] = 30
    cfg["retries"] = 3
    cfg["verbose"] = false

    // Apply overrides
    for entry in overrides.entries() {
        cfg[entry[0]] = entry[1]
    }

    freeze(cfg)
    return cfg
}

flux opts = Map::new()
opts["timeout"] = 60
opts["verbose"] = true

fix config = make_config(opts)
print(config.get("timeout"))   // 60
print(config.get("retries"))   // 3
print(phase_of(config))        // crystal
\end{lstlisting}

\subsection{The Fluid Working Copy}\label{subsec:fluid-copy}
\index{patterns!fluid working copy}

When you need to modify frozen data, \lstinline|thaw()| gives you a working copy:

\begin{lstlisting}[language=Lattice, caption={Working with a thawed copy}]
fix original_scores = freeze([85, 92, 78, 95, 88])

// Create a working copy to add a new score
flux updated = thaw(original_scores)
updated.push(91)
updated = updated.sort()
freeze(updated)

print(original_scores)  // [85, 92, 78, 95, 88] (unchanged)
print(updated)          // [78, 85, 88, 91, 92, 95]
\end{lstlisting}

This immutable-data-with-copies pattern will feel familiar if you've worked with functional programming languages or immutable data structures in Clojure or Elm.

%% ───────────────────────────────────────────────────────────
\section{Exercises}\label{sec:ch04-exercises}

\begin{enumerate}
  \item \textbf{Phase Detective.}
  In the REPL, create variables with \lstinline|let|, \lstinline|flux|, and \lstinline|fix|.
  Use \lstinline|phase_of()| to check the phase of each.
  Freeze a \lstinline|flux| variable and check again.
  Thaw it and check once more.
  Write down what you observe.

  \item \textbf{Strict Mode Workout.}
  Create a file with \lstinline|#mode strict| at the top.
  Try using \lstinline|let| to declare a variable.
  Try assigning to a \lstinline|fix| binding.
  Try freezing a value and then accessing the original name.
  Record the error messages you see.

  \item \textbf{Config Builder.}
  Write a function that builds a database configuration (host, port, database name, connection pool size) using a \lstinline|forge| block.
  The function should accept a map of overrides and return a crystal map with defaults merged in.

  \item \textbf{Freeze Guard.}
  Write a function \lstinline|fn ensure_frozen(value: any) -> any| that returns the value unchanged if it's already crystal, or freezes it if it's fluid.
  Use \lstinline|phase_of()| to decide.

  \item \textbf{The Thaw Game.}
  Start with \lstinline|fix data = freeze([10, 20, 30])|.
  Thaw it, add elements, and freeze again.
  Repeat this three times, building a longer and longer array.
  Verify that each intermediate frozen version is independent.
\end{enumerate}

%% ───────────────────────────────────────────────────────────
\section*{What's Next}\label{sec:ch04-whats-next}

You now understand how Lattice thinks about variables and mutability.
The \lstinline|let|/\lstinline|flux|/\lstinline|fix| keywords give you explicit control over the lifecycle of your data, and the phase system enforces that control at runtime.
We've barely scratched the surface of what the phase system can do---bonds, reactions, sublimation, pressure, and alloys await in Part~III\@.

But first, we need to learn how to \emph{direct} the flow of our programs.
In the next chapter, we'll explore Lattice's control flow constructs: \lstinline|if|/\lstinline|else|, \lstinline|while|, \lstinline|for|, \lstinline|loop|, and the powerful \lstinline|match| expression.
These are the tools that turn sequences of values into programs that make decisions.
